% CODEX: wstaw bezpośrednio przed \section{Krzywizna Minkowskiego i modelu FLRW}
% w 14-krzywizna.tex. Nie wstawiaj ponownie dowodu Ric=0 => R=0 w wymiarze 3.
\subsection*{Krzywizna metryk jednorodnych Nil i Sol}
\label{sec:thur-dop-krzywizna-nil-sol}
Poniższe obliczenia stosują definicję tensora $R$ z
\ref{def:tensor-krzywizny}, tę samą co we wcześniejszych
przykładach, lecz w ramie nieholonomicznej. Wyznaczanie
wszystkich symboli Christoffela w zwykłych współrzędnych
nie jest potrzebne.

\begin{przyklad}[Krzywizna grupy Heisenberga]
\label{prz:thur-dop-krzywizna-nil}
Dla ramy $E_1,E_2,E_3$ z Przykładu~\ref{prz:thur-dop-koneksja-nil}
korzystamy z wyznaczonej tam koneksji i z $[E_1,E_2]=E_3$.
Na przykład
\[
\begin{aligned}
 R(E_1,E_2)E_2
 &=\nabla_{E_1}\nabla_{E_2}E_2
   -\nabla_{E_2}\nabla_{E_1}E_2
   -\nabla_{[E_1,E_2]}E_2\\
 &=0-\tfrac12\nabla_{E_2}E_3-\nabla_{E_3}E_2
 =-\tfrac34E_1.
\end{aligned}
\]
Podobnie $R(E_1,E_3)E_3=\tfrac14 E_1$ oraz
$R(E_2,E_3)E_3=\tfrac14 E_2$. Otrzymujemy
\[
 K_{12}=-\tfrac34,\qquad K_{13}=K_{23}=\tfrac14,
 \qquad \operatorname{Ric}
 =\operatorname{diag}(-\tfrac12,-\tfrac12,\tfrac12),
 \quad \operatorname{Scal}=-\tfrac12.
\]
Przy obliczeniu tensora Ricciego użyliśmy symetrii krzywizny i faktu,
że pozostałe mieszane składowe Ricciego znikają;
można to sprawdzić tym samym wzorem dla $R(E_i,E_j)E_k$.
Geometria Nil ma oba znaki krzywizny sekcyjnej. Dodatnia
wartość Ricciego w kierunku pionowym jest sumą dwóch
dodatnich krzywizn sekcyjnych płaszczyzn zawierających ten kierunek;
w kierunkach poziomych przeważa krzywizna ujemna.
\end{przyklad}

\begin{przyklad}[Krzywizna Sol]
\label{prz:thur-dop-krzywizna-sol}
Dla metryki $g_S$ bierzemy ortonormalną ramę
$E_1=e^t\partial_x$, $E_2=e^{-t}\partial_y$,
$E_3=\partial_t$. Nawiasy mają postać
$[E_3,E_1]=E_1$, $[E_3,E_2]=-E_2$,
$[E_1,E_2]=0$. Wzór~\eqref{eq:thur-dop-koszul-rama}
daje niezerowe pochodne
\[
\nabla_{E_1}E_1=E_3,\quad \nabla_{E_1}E_3=-E_1,
\qquad \nabla_{E_2}E_2=-E_3,\quad
\nabla_{E_2}E_3=E_2;
\]
wszystkie pozostałe $\nabla_{E_i}E_j$ są zerowe.
Dla przykładu
$R(E_1,E_2)E_2=\nabla_{E_1}(-E_3)=E_1$,
a $R(E_1,E_3)E_3=\nabla_{E_1}E_3=-E_1$
(po uwzględnieniu wyrazu z nawiasem $[E_1,E_3]=-E_1$).
Analogicznie otrzymujemy
\[
 K_{12}=1,\qquad K_{13}=K_{23}=-1,
 \qquad \operatorname{Ric}=\operatorname{diag}(0,0,-2),
 \quad \operatorname{Scal}=-2.
\]
W szczególności $\operatorname{Ric}(E_1,E_1)=0$,
choć krzywizny płaszczyzn zawierających $E_1$
są odpowiednio $1$ oraz $-1$. To konkretny
przykład, dlaczego zerowanie pewnej wartości kwadratowej
Ricciego nie jest równoznaczne z płaskością.
\end{przyklad}
